\subsection{トピックと参考文献の対応}
この表は、SWEBOK 第11章の「トピック対参考資料の行列」を、学習用に簡略化したものである。詳しく学ぶときは、各トピックに対応する文献を読む。

\par\smallskip\noindent\refstepcounter{table}\label{tab:ch11-10}表 \thetable: トピックと参考文献の対応の整理\par\smallskip
表 \thetable は、トピックと参考文献の対応の整理を比較・確認しやすい形で整理したものである。\par\smallskip
\begin{longtable}{>{\raggedright\arraybackslash}p{0.42\textwidth}>{\raggedright\arraybackslash}p{0.53\textwidth}}
\toprule
トピック & 主な参考資料 \\
\midrule
モデリング原則 & Budgen; Mellor and Balcer; Sommerville \\
モデルの性質と表現 & Budgen; Sommerville \\
構文・意味論・語用論 & Mellor and Balcer; Sommerville \\
事前条件・事後条件・不変条件 & Mellor and Balcer; Page-Jones \\
構造モデリング & Budgen; Sommerville; Page-Jones; ISO/IEC 19505-1 \\
振る舞いモデリング & Budgen; Mellor and Balcer; Sommerville; Selic and Gerard \\
完全性・整合性の分析 & Sommerville; Wing \\
正しさの分析 & Wing \\
追跡可能性の分析 & Sommerville \\
相互作用の分析 & Mellor and Balcer; Sommerville; Page-Jones \\
経験則的方法 & Budgen; Sommerville; Selic and Gerard; Brambilla et al.; Aarenstrup; Larman; Kiczales et al. \\
形式手法 & Budgen; Sommerville; Wing; Jackson \\
プロトタイピング方法 & Budgen; Sommerville; Brookshear \\
アジャイル方法 & Sommerville; Brookshear; Boehm and Turner; Poppendieck; Ohno; Anderson; Goodpasture; ISO/IEC/IEEE 32675 \\
\bottomrule
\end{longtable}
